% ===========================================================================
% ΠΕΡΙΕΧΟΜΕΝΑ ΒΙΒΛΙΟΥ / ΣΕΙΡΑ ΚΕΦΑΛΑΙΩΝ
% ===========================================================================
% \frontmatter: εισαγωγικές σελίδες, χωρίς κανονική αρίθμηση κεφαλαίων.
% \BookCover: δοκιμαστικό εξώφυλλο.
% \tableofcontents: αυτόματος πίνακας περιεχομένων (θέλει δεύτερη compile).
% \mainmatter: εδώ ξεκινούν τα κεφάλαια με αρίθμηση 1, 2, 3...
% \include{chapters/όνομα}: εισάγει το αντίστοιχο αρχείο .tex.
% Για καινούριο κεφάλαιο δημιούργησε chapters/chapter-02.tex και πρόσθεσε
% \include{chapters/chapter-02}, στη σωστή σειρά, κάτω από τις προηγούμενες.
% ΜΗΝ βάζεις την επέκταση .tex μέσα στο \include.
% ===========================================================================
% Include the real chapters/units here. Both screen and print builds use this file.
\frontmatter
\BookCover
\tableofcontents
\mainmatter
% ===========================================================================
% ΔΟΚΙΜΑΣΤΙΚΟ ΚΕΦΑΛΑΙΟ ΦΥΣΙΚΗΣ — ΟΧΙ Η ΝΕΑ ΣΧΟΛΙΚΗ ΥΛΗ
% ===========================================================================
% Το αρχείο αποτελεί «κατάλογο» των δομικών στοιχείων του template.
% Τα κείμενα, οι τιμές και οι ασκήσεις είναι παραδείγματα δοκιμής.
% Το \BookChapterPhoto παίρνει: τίτλο, υπότιτλο, εικόνα, πηγή, \item στόχους.
% Μπορείς να αλλάξεις το κείμενο χωρίς να αγγίξεις τα κουτιά/γραμματοσειρές.
% ===========================================================================
\BookChapterPhoto{Κίνηση σε μία διάσταση}{Φυσική — δείγμα μακέτας, όχι κείμενο προς δημοσίευση}{assets/chapter-photo-sample.png}{Ενδεικτική εικόνα από το προηγούμενο βιβλίο· προσωρινή για τη μακέτα}{%
\item Να διακρίνεις τη θέση, τη μετατόπιση και το διάστημα.
\item Να επιλέγεις τη θετική φορά ενός άξονα αναφοράς.
\item Να υπολογίζεις τη μετατόπιση από δύο θέσεις.
\item Να ερμηνεύεις το πρόσημο της ταχύτητας.
\item Να διαβάζεις απλά διαγράμματα κίνησης.
}
% Κύρια ενότητα: η μπλε λωρίδα ορίζεται στο universalbook.sty.
\BookSection{Θεωρία}
\section{Περιγραφή της κίνησης}
Για την περιγραφή της ευθύγραμμης κίνησης ορίζουμε έναν άξονα αναφοράς και θετική φορά. Η θέση ενός σώματος προσδιορίζεται από την αλγεβρική τιμή της συντεταγμένης του.
% Τα παλαιότερα δοκιμαστικά σχέδια χρησιμοποιούν πλέον book/axis,
% book/velocity και book/plotA: κοινά στυλ από styles/book-figures.tex.
% KeyFormula = πλαίσιο ανάδειξης βασικής εξίσωσης.
% Εντός του χρησιμοποιούμε κανονική μαθηματική σύνταξη \(...\) ή \[...\].
\begin{KeyFormula}
\(\Delta x=x_{\mathrm{τελ}}-x_{\mathrm{αρχ}}\qquad\text{και}\qquad \upsilon=\dfrac{\Delta x}{\Delta t}\)
\end{KeyFormula}
\begin{center}
\begin{tikzpicture}[x=.95cm,y=.8cm,>=Stealth]
\draw[book/axis] (-5,0)--(5,0) node[right] {$x$};
\foreach \x in {-4,-2,0,2,4} {\draw[BookPrimary] (\x,.10)--(\x,-.10) node[below=2pt]{\footnotesize $\x$};}
\fill[BookPrimary] (-3,0) circle (1.6pt) node[above]{\scriptsize A};
\fill[BookPrimary] (3,0) circle (1.6pt) node[above]{\scriptsize B};
\draw[book/velocity] (-3,.64)--(3,.64) node[midway,above] {\footnotesize $\Delta x$};
\end{tikzpicture}
\end{center}
% «▲ ΠΡΟΣΟΧΗ»: έχει δικό του semantic χρώμα/σύμβολο, άρα κάθε
% προειδοποίηση στο βιβλίο μένει οπτικά πανομοιότυπη.
\begin{AttentionBox}
Η μετατόπιση είναι διανυσματικό μέγεθος, ενώ το διάστημα είναι μονόμετρο. Δεν ταυτίζονται σε κάθε κίνηση.
\end{AttentionBox}
\section{Από το διάγραμμα στη φυσική ερμηνεία}
Η μορφή της γραφικής παράστασης παρέχει πληροφορίες για την κινητική κατάσταση. Στο ακόλουθο δοκιμαστικό διάγραμμα η θέση αυξάνεται γραμμικά με τον χρόνο.
\begin{center}
\begin{tikzpicture}[x=.65cm,y=.40cm,>=Stealth]
\draw[book/axis] (0,0)--(8.0,0) node[right] {$t\,(s)$};
\draw[book/axis] (0,0)--(0,7.0) node[above] {$x\,(m)$};
\foreach \x in {2,4,6}{\draw[BookPrimary] (\x,.10)--(\x,-.10) node[below]{\scriptsize $\x$};}
\foreach \y in {2,4,6}{\draw[BookPrimary] (.10,\y)--(-.10,\y) node[left]{\scriptsize $\y$};}
\draw[book/plotA] (0,1)--(7.2,6.4);
\draw[book/guide] (4,0)--(4,4);
\fill[BookFigPlotAColor] (4,4) circle (2pt);
\end{tikzpicture}
\end{center}
% MethodBox = μεθοδολογία ή αλληλουχία βημάτων επίλυσης.
\begin{MethodBox}
\textbf{Βήμα 1.} Διαβάζουμε τις μονάδες των αξόνων.\par
\textbf{Βήμα 2.} Επιλέγουμε δύο σημεία της γραφικής παράστασης.\par
\textbf{Βήμα 3.} Υπολογίζουμε τον λόγο μεταβολής θέσης προς χρόνο.\par
\textbf{Βήμα 4.} Ελέγχουμε το πρόσημο και τη φυσική σημασία.
\end{MethodBox}

% figure = πλωτό σχήμα, που μπορεί να μετακινηθεί ελαφρά από το LaTeX.
% \caption δίνει μπλε/italic λεζάντα, \label/\ref ασφαλή αρίθμηση.
% htbp = εδώ / πάνω / κάτω / ξεχωριστή σελίδα κατά περίπτωση.
\begin{figure}[htbp]
\centering
\begin{tikzpicture}[x=.86cm,y=.69cm,>=Stealth]
\draw[book/axis] (-5,0)--(5,0) node[right] {$x$};
\foreach \x in {-4,-2,0,2,4}{\draw[BookPrimary] (\x,.07)--(\x,-.07) node[below=2pt]{\scriptsize $\x$};}
\fill[BookPrimary] (-3,0) circle (1.5pt) node[below=5pt]{\scriptsize A};
\fill[BookPrimary] (3,0) circle (1.5pt) node[below=5pt]{\scriptsize B};
\draw[book/velocity] (-3,.7)--(3,.7) node[midway,above=1pt] {$\Delta x$};
\end{tikzpicture}
\caption{Μετατόπιση κινητού από την αρχική θέση Α στην τελική θέση Β.}\label{fig:displacement}
\end{figure}
Όπως φαίνεται στο Σχήμα~\ref{fig:displacement}, η μετατόπιση παριστάνεται με διάνυσμα από την αρχική προς την τελική θέση.
\BookMinorHeading{Πίνακας με επιστημονικά δεδομένα}
\begin{table}[htbp]
\centering
% BookStripedTable: κοινή σκούρα κεφαλίδα, αχνές εναλλάξ γραμμές.
% tabularx: το X προσαρμόζεται αυτόματα στο διαθέσιμο πλάτος.
\begin{BookStripedTable}
\begin{tabularx}{\linewidth}{@{}>{\raggedright\arraybackslash}X>{\centering\arraybackslash}p{25mm}>{\raggedright\arraybackslash}p{35mm}@{}}
\BookTableHeaderRow
\BookTableHead{Μέγεθος} & \BookTableHead{Σύμβολο} & \BookTableHead{Μονάδα S.I.}\\
Μήκος & $\ell$ & μέτρο (m)\\
Μάζα & $m$ & χιλιόγραμμο (kg)\\
Χρόνος & $t$ & δευτερόλεπτο (s)\\
Ταχύτητα & $\upsilon$ & μέτρο ανά δευτερόλεπτο (m/s)\\
Επιτάχυνση & $a$ & μέτρο ανά δευτερόλεπτο$^2$ (m/s$^2$)\\
\end{tabularx}
\end{BookStripedTable}
\caption{Ενδεικτικά φυσικά μεγέθη και μονάδες μέτρησης.}\label{tab:quantities}
\end{table}

% «Λίγη ιστορία»: προαιρετικός τίτλος σε [ ] και κείμενο μέσα.
\begin{HistoryBox}[Η γέννηση της μέτρησης]
Η ιστορία των μετρήσεων στη Φυσική μπορεί να αξιοποιείται ως σύντομο διάλειμμα από τον κύριο κορμό της θεωρίας. Το ένθετο υποστηρίζει εικόνα, λεζάντα και σαφή αναφορά πηγής. Στην τελική έκδοση εδώ θα μπει τεκμηριωμένο ιστορικό κείμενο.
\par\medskip{\footnotesize\itshape Πηγή: συμπληρώνεται μετά τον επιστημονικό έλεγχο.}
\end{HistoryBox}
% Ρήση χωρίς πλαίσιο ή σκίαση. Στο δοκίμιο το κείμενο είναι ΠΡΩΤΟΤΥΠΟ,
% όχι πραγματικό απόφθεγμα ιστορικού προσώπου.
\BookQuote{Η παρατήρηση γεννά ερωτήσεις· η μέτρηση μας βοηθά να τις ελέγξουμε.}
  {Δοκιμαστική ρήση}[Πρωτότυπο κείμενο επίδειξης μορφοποίησης]
\clearpage
% v0.19: κανονική αριθμημένη υποενότητα (εμφανίζεται στον πίνακα περιεχομένων).
% Περιλαμβάνει τα λυμένα, τις ερωτήσεις θεωρίας και τις άλυτες ασκήσεις.
% Δεν γράφουμε χειροκίνητες αποστάσεις ανάμεσα στον τίτλο και την μπλε λωρίδα:
% τις καθορίζει κεντρικά το styles/universalbook.sty.
\section{Εφαρμογές και εξάσκηση}
% Η ίδια κύρια μπλε λωρίδα χρησιμοποιείται και εδώ.
\BookSection{Λυμένα παραδείγματα}
% workedexample αυξάνει μόνο του την αρίθμηση παραδειγμάτων.
% Ο τίτλος μπαίνει στο ανοιχτό μπλε ProblemBox, πάνω στη σκούρα λωρίδα.
\begin{workedexample}{Υπολογισμός μετατόπισης}
% ProblemBox = γαλάζιο πλαίσιο εκφώνησης (όχι λύσης).
\begin{ProblemBox}
Ένα σώμα κινείται πάνω σε ευθεία από τη θέση \(x_1=-3\,\mathrm{m}\) στη θέση \(x_2=5\,\mathrm{m}\).\par
\begin{enumerate}[label=\BookGreek*),font=\color{BookSecondary}\bfseries,leftmargin=7mm,topsep=1pt,itemsep=0pt]
\item Να υπολογιστεί η μετατόπιση.
\item Να προσδιοριστεί η φορά της.
\end{enumerate}
\end{ProblemBox}
% «◉ Πώς σκέφτομαι»: σύντομη στρατηγική, όχι πλήρης λύση.
\begin{ThinkBox}
Ξεκινώ από τον ορισμό της μετατόπισης και εξετάζω το πρόσημο του αποτελέσματος. Το θετικό πρόσημο δηλώνει κατεύθυνση προς τη θετική φορά του άξονα.
\end{ThinkBox}
% Ετικέτα «Λύση» με σκόπιμα ήπια γραφή.
\SolutionHeading
{\color{BookSecondary}\bfseries α)} Από τον ορισμό της μετατόπισης:
\[\Delta x=x_2-x_1=5-(-3)=8\,\mathrm{m}.\]
{\color{BookSecondary}\bfseries β)} Η μετατόπιση έχει θετική αλγεβρική τιμή, επομένως κατευθύνεται προς τη θετική φορά του άξονα.
\end{workedexample}
\begin{workedexample}{Σταθερή ταχύτητα}
\begin{ProblemBox}
Ένα κινητό διανύει μετατόπιση \(\Delta x=60\,\mathrm{m}\) σε χρονικό διάστημα \(\Delta t=12\,\mathrm{s}\). Να υπολογιστεί η αλγεβρική τιμή της σταθερής ταχύτητας.
\end{ProblemBox}
\begin{ThinkBox}
Στην ευθύγραμμη ομαλή κίνηση χρησιμοποιώ τον λόγο \(\Delta x/\Delta t\) και ελέγχω τις μονάδες του αποτελέσματος.
\end{ThinkBox}
\SolutionHeading
\[\upsilon=\frac{\Delta x}{\Delta t}=\frac{60\,\mathrm{m}}{12\,\mathrm{s}}=5\,\mathrm{m/s}.\]
\end{workedexample}
\clearpage
% BookTheorySection = εναλλακτική ονομασία του BookSection.
\BookTheorySection{Ερωτήσεις θεωρίας}
% Δίστηλες σύντομες ερωτήσεις. Για μακροσκελείς ασκήσεις χρησιμοποίησε 1 στήλη.
\begin{multicols}{2}\RaggedRight\raggedcolumns
% Υποκατηγορία με διακριτικό μικρό τίτλο/γραμμή.
% Κάθε \TrueFalseItem προσθέτει αυτόματα και τα δύο αδιάσπαστα κουτάκια.
\BookQuestionHeading{Σωστό / Λάθος}
\begin{enumerate}[label=\arabic*.,leftmargin=13pt]
\TrueFalseItem{Το διάστημα είναι διανυσματικό μέγεθος.}
\TrueFalseItem{Η μετατόπιση εξαρτάται από την αρχική και την τελική θέση.}
\TrueFalseItem{Η ταχύτητα στην ΕΟΚ είναι σταθερή.}
\end{enumerate}
\BookQuestionHeading{Πολλαπλής επιλογής}
Ποια μονάδα χρησιμοποιούμε στο S.I. για την ταχύτητα;
\begin{enumerate}[label=\BookGreek*),font=\color{BookSecondary}\bfseries,leftmargin=16pt]
\item \CheckSquare\hspace{1.1mm}m
\item \CheckSquare\hspace{1.1mm}s
\item \CheckSquare\hspace{1.1mm}m/s
\item \CheckSquare\hspace{1.1mm}m/s$^2$
\end{enumerate}
\columnbreak
\BookQuestionHeading{Ερωτήσεις κρίσεως}
\begin{enumerate}[label=\arabic*.,leftmargin=13pt]
\item Είναι δυνατόν ένα κινητό να έχει μηδενική μετατόπιση και μη μηδενικό διάστημα; Αιτιολόγησε.
\item Πότε το μέτρο της μετατόπισης ισούται με το διάστημα;
\end{enumerate}
\BookQuestionHeading{Διαδοχικά βήματα (i–iv)}
\begin{enumerate}[label=\roman*),leftmargin=15pt,itemsep=0pt]
\item Επιλογή άξονα.
\item Σχεδίαση δυνάμεων.
\item Εφαρμογή νόμου.
\item Έλεγχος μονάδων.
\end{enumerate}
\BookQuestionHeading{Συμπλήρωση κενών}
Συμπληρώστε: Το σύμβολο \(\Delta x\) δηλώνει τη \underline{\hspace{15mm}}.
% Σύντομη απάντηση/υπόδειξη. Κρύβεται μέσω \ShowAnswersfalse.
\Answer{μετατόπιση}
\end{multicols}
% Οι άλυτες είναι «1 | Εκφώνηση», χωρίς ξεχωριστό τίτλο άσκησης.
\BookSection{Άλυτες ασκήσεις}
% [PHY-A-001] = ID στον πηγαίο κώδικα και \label{ex:PHY-A-001}.
% Το {} δηλώνει ότι δεν θέλουμε ορατό τίτλο δίπλα στον αριθμό.
\begin{exercise}[PHY-A-001]{}
Ένα κινητό βρίσκεται αρχικά στη θέση \(x_A=2\,\mathrm m\) και στη συνέχεια στη θέση \(x_B=-7\,\mathrm m\). Να υπολογίσετε τη μετατόπισή του και να περιγράψετε τη φορά.
\Answer{$\Delta x=-9\,\mathrm m$, προς την αρνητική φορά.}
\end{exercise}
\begin{exercise}[PHY-A-002]{}
Ένα σώμα κινείται με σταθερή ταχύτητα. Να εξηγήσετε ποια είναι η επιτάχυνσή του και να σχεδιάσετε ποιοτικά το διάγραμμα επιτάχυνσης–χρόνου.
\Answer{$\alpha=0$.}
\end{exercise}
\begin{exercise}[PHY-A-003]{}
Ένα κινητό κινείται σε ευθεία με σταθερή θετική ταχύτητα. Να σχεδιάσετε τα διαγράμματα $x-t$ και $\upsilon-t$ και να εξηγήσετε τη σημασία της κλίσης τους.
\end{exercise}
\begin{exercise}[PHY-A-004]{}
Ένα σώμα κινείται από τη θέση $x_A=0$ στη θέση $x_B=6\,\mathrm m$ και επιστρέφει στη θέση $x_C=2\,\mathrm m$.
\begin{enumerate}[label=\BookGreek*),leftmargin=18pt,itemsep=0pt]
\item Να βρείτε τη συνολική μετατόπιση.
\item Να υπολογίσετε το συνολικό διάστημα.
\item Να εξετάσετε αν τα δύο μεγέθη ταυτίζονται.
\end{enumerate}
\Answer{$\Delta x=2\,\mathrm m$, $s=10\,\mathrm m$.}
\end{exercise}
\begin{exercise}[PHY-A-005]{}
Δύο σώματα κινούνται πάνω στον ίδιο άξονα. Το πρώτο έχει $\upsilon_1=3\,\mathrm{m/s}$ και το δεύτερο $\upsilon_2=-2\,\mathrm{m/s}$. Να σχολιάσετε τη φορά τους.
\end{exercise}
\begin{exercise}[PHY-A-006]{}
Ένα κινητό επιστρέφει στο σημείο εκκίνησης. Είναι μηδέν η μετατόπιση; Είναι απαραίτητα μηδέν και το διάστημα; Αιτιολογήστε.
\end{exercise}


% Τα Φύλλα Εργασίας και τα Κριτήρια της κίνησης ανήκουν στην Ενότητα 1.
% Γι' αυτό μπαίνουν ΠΡΙΝ από την έναρξη της επόμενης ενότητας.
% ===========================================================================
% ΔΟΚΙΜΑΣΤΙΚΟ ΦΥΛΛΟ ΕΡΓΑΣΙΑΣ
% ===========================================================================
% \WorksheetOpening: ξεχωριστή μπλε επικεφαλίδα και αύξων αριθμός φύλλου.
% \WorksheetSection: μπλε μπάρα για επιμέρους δραστηριότητα.
% \CheckSquare: τετραγωνάκι επιλογής, με μέγεθος από book-config.tex.
% \AnswerLine/\AnswerGrid: χώρος για γραπτές απαντήσεις μαθητή.
% ===========================================================================
% The worksheet is intentionally independent of the subject text.
\WorksheetOpening{ΚΙΝΗΜΑΤΙΚΗ}{Θέση, μετατόπιση και ερμηνεία διαγραμμάτων}
\begin{LearningGoals}
\begin{enumerate}
\item Να υπολογίζουν μετατόπιση από δύο θέσεις.
\item Να επιλέγουν σωστή μονάδα φυσικού μεγέθους.
\item Να αιτιολογούν σύντομα τις απαντήσεις τους.
\end{enumerate}
\end{LearningGoals}
\WorksheetSection{Α. Ελέγχω την κατανόησή μου}
% ΣΤΥΛ v0.16:
% Ο τίτλος «Σωστό ή λάθος;» έχει μόνο μαύρα έντονα γράμματα.
% Το υπόλοιπο κείμενο είναι κανονική μαύρη Inter (όχι μπλε ή bold).
% Η μορφοποίηση γίνεται κεντρικά από το exercise στο universalbook.sty.
\begin{exercise}[]{Σωστό ή λάθος;}
Η μετατόπιση ενός κινητού είναι πάντα ίση με το διάστημα που διήνυσε.\hspace{.5em}\hfill\TrueFalseChoices
\end{exercise}
\AnswerLine
\WorksheetSection{Β. Εφαρμόζω}
\begin{exercise}[]{Υπολογισμός μετατόπισης}
Ένα κινητό μετακινείται από τη θέση \(x_A=-2\,\mathrm m\) στη θέση \(x_B=6\,\mathrm m\). Να υπολογίσετε τη μετατόπιση και να σχολιάσετε το πρόσημό της.
\end{exercise}
\AnswerGrid{2}
\WorksheetSection{Γ. Εξηγώ}
\begin{exercise}[]{Ερώτηση κρίσεως}
Να περιγράψετε μια κίνηση στην οποία το διάστημα είναι θετικό, αλλά η συνολική μετατόπιση είναι μηδέν.
\end{exercise}
\AnswerLine
\vfill\begin{center}{\footnotesize\color{BookSecondary}ΔΕΙΓΜΑ ΦΥΛΛΟΥ ΕΡΓΑΣΙΑΣ · ΔΕΝ ΑΠΟΤΕΛΕΙ ΤΕΛΙΚΗ ΔΙΔΑΚΤΙΚΗ ΥΛΗ}\end{center}

% =======================================================================
% ΚΡΙΤΗΡΙΟ ΑΞΙΟΛΟΓΗΣΗΣ — ΔΟΚΙΜΑΣΤΙΚΟ ΠΕΡΙΕΧΟΜΕΝΟ
% =======================================================================
% Τα συγκεκριμένα διαγωνίσματα αφορούν την Ενότητα 1 (Κινηματική).
% Τοποθετούνται πριν από το επόμενο κεφάλαιο στο book-content.tex,
% ώστε να εμφανίζονται ΕΝΤΟΣ της Ενότητας 1 στα περιεχόμενα.
% Ολες οι μορφοποιήσεις ορίζονται κεντρικά.
% Τα παρακάτω ΔΕΝ αποτελούν ολοκληρωμένο επιστημονικά διαγώνισμα προς χρήση.
% Είναι σχεδιαστικές δοκιμές των περιβαλλόντων και των σημείων.
\begin{Assessment}[student=true,date=true,duration={90 λεπτά}]{Ευθύγραμμη κίνηση}
\AssessmentTopic{Α}
\begin{AssessmentQuestion}[5]\label{asq:KRI1-A1}
Η ταχύτητα ενός σώματος παραμένει σταθερή κατά μέτρο και κατεύθυνση.
Η επιτάχυνσή του είναι:
\begin{enumerate}[label=\BookGreek*. ,leftmargin=7mm,itemsep=1pt]
\item θετική. \item αρνητική. \item μηδενική. \item μεταβαλλόμενη.
\end{enumerate}
\end{AssessmentQuestion}
\begin{AssessmentQuestion}[5]
Να χαρακτηρίσετε την πρόταση ως σωστή ή λανθασμένη:
«Το διάστημα είναι πάντοτε ίσο με το μέτρο της μετατόπισης».
\TrueFalseChoices
\end{AssessmentQuestion}
\begin{AssessmentQuestion}[5]
Το μέτρο της μέσης διανυσματικής ταχύτητας ενός σώματος
υπολογίζεται από τη σχέση $|\bar{\upsilon}|=\frac{|\Delta x|}{\Delta t}$.
Να αναφέρετε τις μονάδες SI των μεγεθών $\Delta x$ και $\Delta t$.
\end{AssessmentQuestion}
\begin{AssessmentQuestion}[5]
Ένα κινητό κινείται από τη θέση $x_1=-2\,\mathrm{m}$ στη θέση
$x_2=6\,\mathrm{m}$. Ποια είναι η αλγεβρική μετατόπισή του;
\end{AssessmentQuestion}
\begin{AssessmentQuestion}[5]
Να ορίσετε τη μέση αριθμητική ταχύτητα σε μια κίνηση.
\end{AssessmentQuestion}

\AssessmentTopic{Β}
\begin{AssessmentQuestion}[12]
Δύο κινητά Α και Β κινούνται στον ίδιο άξονα.
Τα διαγράμματα θέσης--χρόνου φαίνονται παρακάτω.
\begin{center}
\begin{tikzpicture}[scale=.67]
 \draw[book/axis] (0,0)--(6.4,0) node[right]{$t\ (\mathrm{s})$};
 \draw[book/axis] (0,-.4)--(0,3.9) node[above]{$x\ (\mathrm{m})$};
 \draw[book/velocity] (0,.35)--(5.4,3.55) node[right]{A};
 \draw[book/accel] (0,3.1)--(5.4,1.05) node[right]{B};
\end{tikzpicture}
\end{center}
Ποιο κινητό έχει θετική και ποιο αρνητική αλγεβρική ταχύτητα;
Να αιτιολογήσετε την επιλογή σας.
\end{AssessmentQuestion}
\begin{AssessmentQuestion}[13]
Δύο σώματα διανύουν σε ίσους χρόνους διαστήματα $12\,\mathrm{m}$ και
$18\,\mathrm{m}$ αντίστοιχα. Ποιο έχει μεγαλύτερη μέση αριθμητική ταχύτητα;
\end{AssessmentQuestion}

\AssessmentTopic{Γ}
\begin{AssessmentQuestion}[10]\label{asq:KRI1-G1}
Ένα κινητό ξεκινά από τη θέση $x_0=2\,\mathrm{m}$ και κινείται με
σταθερή ταχύτητα $\upsilon=3\,\mathrm{m/s}$.
Να γράψετε την εξίσωση της θέσης του.
\end{AssessmentQuestion}
\begin{AssessmentQuestion}[15]
Για το προηγούμενο κινητό, να υπολογίσετε τη θέση του τη χρονική στιγμή
$t=4\,\mathrm{s}$ και τη μετατόπισή του στο ίδιο χρονικό διάστημα.
\end{AssessmentQuestion}

\AssessmentTopic{Δ}
\begin{AssessmentQuestion}[9]
Ένα σώμα αρχίζει να κινείται από την ηρεμία με σταθερή επιτάχυνση
$\alpha=2\,\mathrm{m/s^2}$. Να υπολογίσετε την ταχύτητά του τη στιγμή
$t=5\,\mathrm{s}$.
\end{AssessmentQuestion}
\begin{AssessmentQuestion}[8]
Να υπολογίσετε τη μετατόπισή του στα πρώτα $5\,\mathrm{s}$.
\end{AssessmentQuestion}
\begin{AssessmentQuestion}[8]
Να σχεδιάσετε ποιοτικά το διάγραμμα ταχύτητας--χρόνου για
$0\leq t\leq5\,\mathrm{s}$.
\end{AssessmentQuestion}
% Προαιρετικός έλεγχος στο αρχείο LOG: 25+25+25+25=100.
\AssessmentPointsCheck[100]
\end{Assessment}

% Δεύτερο, σύντομο κριτήριο: ΟΥΤΕ στοιχεία μαθητή ΟΥΤΕ ημερομηνία
% ΟΥΤΕ διάρκεια. Κάποια ερωτήματα δεν έχουν μονάδες.
\begin{Assessment}{Σύντομη επανάληψη: μετατόπιση}
\AssessmentTopic{Α}
\begin{AssessmentQuestion}
Ποια είναι η διαφορά ανάμεσα στη μετατόπιση και στο διάστημα;
\end{AssessmentQuestion}
\begin{AssessmentQuestion}[5]
Δίνεται $x_A=4\,\mathrm{m}$ και $x_B=-1\,\mathrm{m}$.
Να βρείτε τη μετατόπιση από το Α στο Β.
\end{AssessmentQuestion}
\AssessmentTopic{Β}
\begin{AssessmentQuestion}
Να εξηγήσετε αν είναι δυνατόν η μετατόπιση ενός κινητού να είναι μηδέν,
ενώ το διάστημα που έχει διανύσει είναι μεγαλύτερο από μηδέν.
\end{AssessmentQuestion}
\end{Assessment}


% Επόμενο κεφάλαιο, με δικά του φύλλα και κριτήρια όταν προστεθούν.
% ===========================================================================
% ΔΟΚΙΜΑΣΤΙΚΟ ΚΕΦΑΛΑΙΟ ΜΑΘΗΜΑΤΙΚΩΝ
% ===========================================================================
% Η ίδια μακέτα εξυπηρετεί Μαθηματικά και άλλα επιστημονικά βιβλία.
% Δεν απαιτείται διαφορετικό universalbook.sty: το περιεχόμενο αλλάζει.
% Δες παρακάτω τις εντολές definitionbox, theorembox, proof, TikZ και figure.
% ===========================================================================
\BookChapterPhoto{Συναρτήσεις και γραφήματα}{Μαθηματικά — δοκιμή της ίδιας εκδοτικής μακέτας}{assets/math-chapter-art.pdf}{Διανυσματική εικονογράφηση, ενσωματωμένη στο XeLaTeX}{%
\item Να αναγνωρίζεις συναρτήσεις πρώτου βαθμού.
\item Να ερμηνεύεις τους συντελεστές τους.
\item Να συνδέεις τύπους με γραφικές παραστάσεις.
\item Να γράφεις σαφή μαθηματική αιτιολόγηση.
}
\BookSection{Ορισμός και βασική ιδέα}
\section{Η συνάρτηση πρώτου βαθμού}
Ο τύπος \(f(x)=\alpha x+\beta\), με \(\alpha\neq0\), περιγράφει μια συνάρτηση πρώτου βαθμού. Η γραφική της παράσταση είναι ευθεία.
% Ορισμός: αριθμείται αυτόματα ανά κεφάλαιο.
\begin{definitionbox}[Συνάρτηση πρώτου βαθμού]
\(f(x)=\alpha x+\beta\), όπου \(\alpha\neq0\) και \(\alpha,\beta\in\mathbb R\).
\end{definitionbox}

\begin{table}[htbp]
\centering
% v0.23: ο λιτός πίνακας έχει την ΙΔΙΑ σκούρα κεφαλίδα
% και τις ΙΔΙΕΣ αποστάσεις με τον χρωματιστό.
% Μόνη διαφορά: οι γραμμές δεδομένων δεν αποκτούν φόντο.
\begin{BookPlainTable}
\begin{tabularx}{.82\linewidth}{@{}>{\centering\arraybackslash}X>{\centering\arraybackslash}X@{}}
\BookTableHeaderRow
\BookTableHead{$x$} & \BookTableHead{$f(x)=2x+1$}\\
$-1$ & $-1$\\
$0$ & $1$\\
$1$ & $3$\\
\end{tabularx}
\end{BookPlainTable}
\caption{Λιτός πίνακας: ίδιος τίτλος κεφαλίδας, λευκές γραμμές δεδομένων.}
\end{table}

\BookMinorHeading{Παρατήρηση}
Η παράμετρος \(\alpha\) καθορίζει την κλίση της ευθείας, ενώ η παράμετρος \(\beta\) καθορίζει την τεταγμένη του σημείου τομής με τον κατακόρυφο άξονα.
\begin{AttentionBox}
Στο γενικό πρότυπο τα χρώματα σημαίνουν τη λειτουργία ενός πλαισίου, όχι το γνωστικό αντικείμενο. Το ίδιο πλαίσιο «Προσοχή» χρησιμοποιείται με συνέπεια σε Φυσική, Μαθηματικά και άλλα βιβλία.
\end{AttentionBox}
\BookSection{Απόδειξη — υπόδειγμα}
% Θεώρημα: αριθμείται αυτόματα ανά κεφάλαιο.
\begin{theorembox}[Πρόσθεση ίσων ποσοτήτων]
Αν \(a=b\), τότε \(a+c=b+c\) για κάθε πραγματικό αριθμό \(c\).
\end{theorembox}
% Απόδειξη: η ονομασία «Απόδειξη» έρχεται από τον κεντρικό config.
\begin{proof}
Προσθέτοντας τον ίδιο αριθμό στα δύο μέλη της ισότητας, διατηρούμε την ισότητα.
\end{proof}
\begin{workedexample}{Εφαρμογή σε συνάρτηση}
\begin{ProblemBox}
Για τη συνάρτηση \(f(x)=2x+1\), να υπολογιστούν \(f(0)\) και \(f(3)\).
\end{ProblemBox}
\begin{ThinkBox}
Αντικαθιστώ στη μεταβλητή \(x\) την τιμή που μου ζητείται, χωρίς να αλλάζω τον τύπο της συνάρτησης.
\end{ThinkBox}
\SolutionHeading
\[f(0)=2\cdot0+1=1,\qquad f(3)=2\cdot3+1=7.\]
\end{workedexample}
\BookSection{Άσκηση}
\begin{exercise}[MATH-A-001]{}
Να βρεθεί η τιμή του \(x\) για την οποία \(f(x)=9\), όταν \(f(x)=2x+1\).
\Answer{$x=4$.}
\end{exercise}


% Προαιρετικό λυσάρι κριτηρίων, σε ξεχωριστό αρχείο.
% ===========================================================================
% ΕΝΔΕΙΚΤΙΚΕΣ ΛΥΣΕΙΣ ΚΡΙΤΗΡΙΩΝ ΑΞΙΟΛΟΓΗΣΗΣ — ΞΕΧΩΡΙΣΤΟ ΑΡΧΕΙΟ
% ===========================================================================
% Η αναφορά γίνεται με \label μέσα στο αντίστοιχο AssessmentQuestion.
% Δεν αντιγράφουμε ξανά την εκφώνηση. Αν αλλάξει η σειρά ερωτημάτων,
% το \ref ενημερώνεται στη δεύτερη μεταγλώττιση.
% Οι λύσεις εδώ είναι ενδεικτικές και όχι πλήρες λυσάρι του δοκιμίου.
\clearpage
\thispagestyle{plain}
\phantomsection
\addcontentsline{toc}{chapter}{Λύσεις κριτηρίων αξιολόγησης}
\BookRememberRunningChapter{Λύσεις κριτηρίων αξιολόγησης}%
\gdef\BookRunningChapterNumber{}%
\markboth{Λύσεις κριτηρίων}{Λύσεις κριτηρίων}
\BookSection{Λύσεις κριτηρίων αξιολόγησης}
\BookMinorHeading{Κριτήριο 1 — Ευθύγραμμη κίνηση}
\begin{AssessmentSolution}{asq:KRI1-A1}
Σωστή επιλογή: γ. Για σταθερή διανυσματική ταχύτητα, η επιτάχυνση
είναι μηδενική.
\end{AssessmentSolution}
\begin{AssessmentSolution}{asq:KRI1-G1}
Για ευθύγραμμη ομαλή κίνηση η εξίσωση θέσης είναι
\[
  x=x_0+\upsilon t=2+3t\quad (\mathrm{SI}).
\]
\end{AssessmentSolution}

